\chapter{Семантика операторов интерпретации языка Smile}
\label{semantics}

Наличие статической типизации привело, в сравнеии с Пифагором, к пересмотру семантики операций над контейнерными данными. В частности этому содействовало то, что у вектора, роя и потока все данные однотипные, а тип данных известен. Это позволяет выполнять над ними без дополнительной проверки типов не только массовые функции над всеми элементами, но и рассматривать каждый контейнер как единое целое, представляющее активный объект.

Использование статической типизации также привело к разделению оператора интерпретации на два разных вида: одиночный (одноаргументный) и групповой (массовый, поэлементный). Одиночный оператор интерпретации, обозначаемый через <<\verb|:|>> (постфиксная форма) или <<\verb|^|>> (префиксная форма) предназначен для задания обычных функций, воспринимающих аргумент в качестве единого целого. Массовый оператор интерпретации используется для задания вычислений над каждым однотипным элементом контейнера, порождая на выходе контейнер с элементами тип которых соответствует типу результата выполняемой функции. Обозначается двойным значком <<\verb|::|>> для постфиксной или <<\verb|^^|>> для префиксной форм соответственно.

Использование разных обозначений позволяет однозначно применять функцию с одним и тем же именем в разных контекстах. Например, функция вычитания <<\verb|-|>> над аргументом (10, -3), воспринимаемом как вектор, состоящий из двух целых чисел, порождает следующие значения:
\begin{verbatim}
    // двуместная функция вычитания над одним аргументом
    (10, -3):-  =>  13

    // функция смены знака, массово применяемая к двум однотипным аргументам
    (10, -3)::- => (-10, 3)
\end{verbatim}

Разделение оператора интерпретации на одиночный и массовый позволяет ввести более гибкий одноаргументный набор дополнительных функций для потока и роя, обеспечивающих обработку асинхронно поступающих данных.

Одиночный оператор интерпретации обрабатывает как функции, определенные семантикой языка (предопределенные функции), так и функции, разрабатываемые программистом.

\textit{\textbf{Примечание.} Вызывает сомнение фраза:
В отличие от него семантика массового оператора интерпретации ориентирована только на предопределенные функции, обрабатывающие композиции данных.}

\section{Использование функций с одиночным оператором интерпретации}

Одиночный оператор интерпретации по сути является ключевым, так как не ведет к последующей декомпозиции данных, а рассматривает свои аргументы как единое целое, для каждой пары которых определена семантика выполнения. Эта семантика четко определяет соответствие между функцией и данными. Оператор изначально отслеживает, чтобы аргумент-данные соответствовал сигнатуре аргумента-функции, после чего обеспечивает выполнение функции на обработку данных. При отсутствии этого соответствие результатом является ошибка и прерывание программы.

\subsection{Предопределенные функции}

Существует большое число предопределенных функций, сигнатура которых и семантика вычислений определены изначально. Эти функции используются как для выполнения базовых арифметических и логических операций, так и для обеспечения различных композиций и декомпозиций данных, что осуществляется в дополнение к операторам языка.

Предопределенные функции и данные формируются на основе атомов, каждый из которых может быть в роли как аргумента, так и функции оператора интерпретации. При этом ряд атомарных артефактов могут выступать только в роли данных, другие - в роли функций, третьи - в той и другой.

Для идентификации предопределенных функций в основном используются специальные символы. Обычно за ними закрепляются вычислительные операции, традиционные для этих символов в большинстве существующих языков программирования. Задаются допустимые аргументы и значения этих функций. При этом типы аргумента и результата являются именованными типами и определяется аксиоматически при описании функции. Допускается перегрузка предопределенных функций, при которой формируются новые сигнатуры.

\textit{\textbf{Примечание.}
В данный момент я пытаюсь выделить подмножество предопределенных функций, предназначенных для формирования каркаса ФПП программы, характерного для Праязыка. В таком виде программа вряд ли может создаваться, так как минимальный набор манипуляций с данными (сравнения и т.д.) необходим даже для самого базового языка.}

\textit{В рассматриваемой версии не все спецзнаки имеют определенную семантику. Предполагается, что в дальнейшем она будет разрабатываться и уточняться. Наряду со специальными знаками в качестве предопределенных функций могут выступать и идентификаторы, которые, как и ключевые слова, запрещены для других применений.}

Ниже приводится описание семантики предопределенных функций, используемых в текущей версии. Следует отметить, что для аргументов, тип которых при описании не задан результатом будет ошибка: <<неправильное использование функции>>.

Можно выделить группу предопределенных функций общего назначения, обеспечивающих универсальные преобразования данных независимо от их типа и структуры. Они позволяют работать с разными артефактами.

\subsection{Использование знака <<|>>}

Знак \verb||| задает функцию нахождения длины контейнерного типа. Аргументом-данными могут быть вектор, кортеж, рой. Результат – целое число, задающее количество элементов.

Данная функция имеет следующие сигнатуры::

\begin{verbatim}
    :| << func any(*) -> int
    :| << func any[*] -> int
    :| << func (...) -> int
\end{verbatim}

Если аргумент не является одним из указанных контейнерных типов, то результатом является ошибка операции интерпретации \verb|BASEFUNCERROR|. Использование функции позволяет проверить размер аргумента перед обработкой, а лишь затем начать выделение его элементов. Следует при этом отметить, что выделение элементов для кортежа, в соответствии с описанным выше, может осуществляться только если тип этих элементов определен на этапе компиляции или если есть возможность приведения к векторному типу.

Примеры::

\begin{verbatim}
    (a, n, (q, w), s):| => 4
    [1, 2, 3, 4, 5]:| => 5
    ((1, 2, (f, d), x)):| => 1
\end{verbatim}

\subsubsection{Использование знака <<?>>}

Функция <<\verb|?|>> осуществляет единичную селекцию, то есть, ищет в булевом векторе данных позицию единственного истинного значения и выводит его целочисленный номер. Функция полезна для организации единичного выбора. Например:
\begin{verbatim}
    (false, true, false, false):? => 1
\end{verbatim}
Данная функция имеет следующие сигнатуры:
\begin{verbatim}
    :? << func bool(*) -> int
    :? << func bool[*] -> int
\end{verbatim}
Если список булевских величин состоит только из ложных значений, или в том случае, когда в нем имеется несколько истинных значений формируется ошибка базовой функции:
\begin{verbatim}
    (false, false, false):? => BASEFUNCERROR
    (true, false, true, false):? => BASEFUNCERROR
\end{verbatim}

\subsubsection{Использование знака <<??>>}

Функция <<\verb|??|>>, осуществляет вычисление номеров позиций истинных булевых констант в булевом векторе данных. То есть, используется как селектор множества позиций. В качестве результата формируется целочисленный вектор с номерами тех элементов булево вектора, чьи значения были равны \verb|true|. Функция полезна для организации выборочного продолжения дальнейших вычислений с использованием множественного селектора. Полученные целочисленные значения используются для выбора элементов из контейнерных данных. Например:
\begin{verbatim}
    (true, false, true, false, false, true):?? => (0, 2, 5)
\end{verbatim}
Данная функция имеет следующие сигнатуры:
\begin{verbatim}
    :?? << func bool(*) -> int(*)
    :?? << func bool[*] -> int[*]
\end{verbatim}
Если список состоит только из ложных значений, на выходе формируется вектор, не содержащий элементов, автоматически преобразуемый в ошибку базовой функции:
\begin{verbatim}
    (false, false, false):?? => () => BASEFUNCERROR
\end{verbatim}

\subsubsection{Использование знака <<..>>}

Знак <<\verb|..|>> используется в качестве функции, формирующей вектор или рой из числовых атомов. В качестве аргумента может выступать трехэлементный вектор или рой, в котором первое число задает начало интервала, второе - его конец, а третье - шаг. Числа могут быть как целые, так и действительные, но одного типа, а шаг принимать как положительные, так и отрицательные значения. Необходимо отсутствие расхождений между значением шага и границами интервала.

Кроме этого аргумент функции может быть целочисленным двухэлементным вектором или роем. В этом случае первый элемент определяет нижнюю целочисленную границу интервала, а второй верхнюю. При этом нижняя граница должна быть меньше или равной верхней границе, а шаг по умолчания принимается равным единице.

При некорректном задании границ интервала, идентифицированном во время выполнения, данная функция возвращает ошибку \verb|BOUNDERROR|. Если некорректно задана структура контейнера или тип его элементов, то возвращается ошибка периода компиляции.

Данная функция имеет следующие сигнатуры:
\begin{verbatim}
    :.. << func int(2) -> int(*)
    :.. << func double(2) -> double(*)
    :.. << func int(3) -> int(*)
    :.. << func double(3) -> double(*)
    :.. << func int[2] -> int[*]
    :.. << func double[2] -> double[*]
    :.. << func int[3] -> int[*]
    :.. << func double[3] -> double[*]
\end{verbatim}

Примеры:
\begin{verbatim}
    (-3.5, 2.0, 1.5):.. => (-3.5, -2.0, -0.5, 1.0)
    [1, 5]:.. => [1, 2, 3, 4, 5]
    (2, 1):.. => BOUNDERROR
\end{verbatim}

\textit{\textbf{Примечание.} В дальнейшем предполагается введение интервальных констант, задающих необходимый диапазон. Реализация действительных чисел на первом этапе не предполагается.}

\subsubsection{Использование функции <<dup>>}

Функция обеспечивает создания вектора из одинаковых элементов путем дублирования. Аргумент --- кортеж, первый элемент которого --- значение некоторого предопределенного именованного типа, а второй --- положительная целочисленная константа. Результат --- вектор, элементами которого являются копии первого элемента аргумента, а количество элементов равно значению второго элемента аргумента.

Пример:

\begin{verbatim}
    (10, 5):dup => (10, 10, 10, 10, 10)
\end{verbatim}

Данная функция имеет следующую сигнатуру::

\begin{verbatim}
    dup << func (any, int) -> any(*)
\end{verbatim}

При некорректном задании количества дубликатов, идентифицированном во время выполнения, данная функция возвращает ошибку \verb|BOUNDERROR|.

\subsubsection{Использование знака <<\#>>}

Данный знак предполагает преобразование двух кортежей одинаковой длины, размещенных в кортеже в кортеж, состоящий из множества двоек. В качестве исходных аргументов могут выступать как атомарные элементы, так и составные конструкции различного типа. Формируемые пары являются кортежами. Например:

\begin{verbatim}
    ((1,2,3),(4,5,6)):# => ((1,4),(2,5),(3,6))
    (((1,4),(2,5),(3,6)), (+,-,*)):# => (((1,4),+),((2,5),-),((3,6),*))
\end{verbatim}
Данная функция может использоваться в дальнейшем для реализации различных групповых операций над данными. Она имеет следующую сигнатуру:

\begin{verbatim}
    :# << func tuple(tuple(t10, t11,...), tuple(t20, t21...))
          -> tuple(tuple(t10,t20), tuple(t11,t21),...)
\end{verbatim}
При некорректном задании аргументов, идентифицированном во время выполнения, данная функция возвращает ошибку \verb|BOUNDERROR|. При несовпадении размеров кортежей функция возвращает ошибку \verb|SIZEERROR|.

\subsection{Использование данных}

Данные тоже могут допускать различное толкование в зависимости от того, на какой из входов оператора интерпретации они поступают.

\subsubsection{Использование в качестве функции сигнальной константы <<!>>}

Использование сигнальной константы в качестве функции ориентировано на отображение аргумента, поступившего на вход данных в качестве результата.
Функция имеет следующую сигнатуру:
\begin{verbatim}
    :! << func any -> any
\end{verbatim}

То есть, в качестве функции знак <<\verb|!|>> интерпретируется как пустая операция, осуществляющая отображение аргумента в качестве результата. Данном знаку можно сопоставить следующее определение функции:
\begin{verbatim}
    :! << func x@any -> any {
        x:return
    }
\end{verbatim}

Если аргумент является вектором, группой или атомом, то происходит его выдача в качестве результата.  Если же аргумент определен как задержка, то предварительно оператор интерпретации осуществляет ее раскрытие, ведущее к вычислению задержанного выражения. Полученное значение передается функции, которая и возвращает его уже без дополнительных изменений.

Пример:
\begin{verbatim}
    {(2,3):+}:! => 5
\end{verbatim}

Сигнальная константа также может использовать как аргумент, то есть, поступать на вход данных. В этом случае она интерпретируется как отсутствие аргумента.  Обычно используется с теми функциями без параметров. То есть она воспринимается как сигнал, запускающий функцию в требуемый момент времени. Например, пусть будет задана функция:
\begin{verbatim}
    sin_pi_div_4 << func -> double {
        (3.14, 4.0):/:sin >>return
    }
\end{verbatim}

Тогда ее вызов будет оформлен следующим образом:
\begin{verbatim}
    !:sin_pi_div_4
\end{verbatim}

\subsubsection{Использование целых чисел в качестве функции}

Целые положительные числа могут использоваться в качестве селекторов для выбора элементов из векторов, кортежей и роев. То есть, они могут выступать в качестве индексов. Получение элемента из кортежа можно показать на следующем примере:
\begin{verbatim}
    (10, 9, 8, 7, 6, 5, 4, 3, 2, 1):5 => 6
\end{verbatim}

\textit{\textbf{Примечание.}
Вариант с отрицательными индексами, существовавший в Пифагоре, убран, так как он ведет к динамическому изменению размеров контейнерных типов, что вряд ли целесообразно при использовании статической типизации. Поэтому корректным считаются только  нулевой и положительный индекс.}

Данная функция имеет следующую сигнатуры::

:@int << func any(*) -> any
:@int << func (...) -> any
:@int << func any[*] -> any
Знак типа "@" указывает на то что подставляется не имя  \verb|int|, а некоторое целочисленное значение, которое может быть константой или получено в ходе вычислений. Момент формирования значения в контейнерном объекте связан с его типом.

При использовании индекса с кортежем необходимо, чтобы все его данные кортежа были одного типа, что позволяет получить на выходе предопределенный тип данных неявно приводимый к вектору. Кроме этого возможна ситуация, когда разные типы элементов кортежа известны во время компиляции, а значением селектора является константное выражение, что позволяет вывести тип результата. Несоответствие этому условию определяется во время компиляции. Если в качестве данных выступает рой, то оператор интерпретации запускается на выполнения не дожидаясь поступления всех элементов роя. Достаточно наличия целочисленного значения, определяющего функцию селекцию и прихода в рой элемента с указанным индексом. То есть в данной ситуации осуществляется асинхронный выбор элемента, что может ускорить вычисления, но ведет к дополнительным расходам при реализации механизма.

В случае отрицательного значения числа порождается ошибка \verb|BOUNDERROR|, сигнализирующая о неправильном индексе.

\subsubsection{Использование булевых значений в качестве функций}

Функция, определяемая булевым значением, осуществляет выбор первого или второго элемента двухэлементного вектора в зависимости от значения. Значение <<\verb|true|>> обеспечивает выбор первого элемента вектора, а значение <<\verb|false|>> --- второго. Оба значения должны принадлежать одному и тому же типу. По сути это селектор значения, реализующий простейшее условие выбора. В качестве аргумента могут выступать двухэлементные: вектор, рой или кортеж с элементами одного типа, приводимый к двухэлементному вектору. Отличие между вектором (кортежем) и роем определяется в подходе к синхронизации вычислений, как и в случае использования в качестве функций целочисленных селекторов.

Функция, определяемая булевским значением, имеет следующие сигнатуры:
\begin{verbatim}
    :@bool << func (any, any) -> any
    :@bool << func any(2) -> any
    :@bool << func any[2] -> any
\end{verbatim}

Примеры:
\begin{verbatim}
    (10, 20):true  => 10
    (10, 20):false => 20
    [10, 20]:true  => 10
    [10, 20]:false => 20
\end{verbatim}

В случае вектора меньшей или большей размерности порождается ошибка \verb|BOUNDERROR|, сигнализирующая о неправильном обращении. Следует также отметить, что порядок выбора не связан с числовыми значениями, используемыми для представления <<\verb|true|>> и <<\verb|false|>>, а определяется в соответствии с общепринятой логикой функционирования условного оператора <<\verb|if...then...else|>> традиционных языков программирования.

\subsection{Функции, выполняющие вычисления}

Ряд специальных знаков зарезервированы для выполнения функции, свойственных их традиционному применению во многих языках программирования. Следует отметить, что в ряде случаев они изначально перегружены в большей степени, чем соответствующие знаки в ряде других языков. Это обусловлено отсутствием преобразований типов по умолчанию, что обеспечивает большую однозначность и понимание того, в каком контексте используется тот или иной знак. Практически все эти знаки имеют тип \verb|none:func|.

\textit{\textbf{Примечание.}
Следует обратить внимание на присутствие в тексте действительных чисел, которые пока реализовывать не планируется. Но они ведь все равно появятся...}

\subsubsection{Сложение <<+>>}

Знак <<\verb|+|>> используется для задания бинарных и унарных операций. В случае бинарных операций его аргументом является двухэлементный кортеж. Интерпретация знака <<\verb|+|>> в качестве функции зависит от типов элементов кортежа. Если аргумент является двухэлементным кортежем числовых атомов (целых или действительных), то выполняется арифметическое сложение. При сложении двух целых чисел результат всегда является целым числом. В этом случае автоматическое преобразование к действительному числу не происходит даже при переполнении. Сложение действительных чисел осуществляется, когда аргументом функции является двухэлементный кортеж действительных чисел. Результатом в этом случае является действительное число. Следует отметить отсутствие неявных преобразований при выполнении сложения. Исключением является преобразование целочисленной константы "0" в действительный нуль "0.0".

Если аргумент является числовым атомом, то он выдается в качестве результата без каких-либо преобразований.

Аргумент функции <<\verb|+|>> может также быть булевым вектором длиной, равной 2. Результатом в этом случае является дизъюнкция (логическое <<или>>). Допускается одноэлементный булев атом, порождающий в качестве результата значение этого же элемента.

Во всех остальных случаях результатом является ошибка операции интерпретации \verb|BASEFUNCERROR|.

Функция, имеет следующие сигнатуры:
\begin{verbatim}
    :+ << func (int, int) -> int
    :+ << func (double, double) -> double
    :+ << func (bool, bool) -> bool
    :+ << func int -> int
    :+ << func double -> double
    :+ << func bool -> bool
\end{verbatim}

Примеры.
\begin{verbatim}
    (3, 5):+ => 8
    (3.0, 5.0):+ => 8.0
    (3, 5.0):+ => BASEFUNCERROR
    (5):+ => BASEFUNCERROR
    5:+ => 5
    (max_integer, 1):+ => 0
    (true, false):+ => true
    true:+ => true
    false:+ => false
    (true):+ => BASEFUNCERROR
\end{verbatim}

\textit{\textbf{Примечание.}
Арифметические операции над действительными числами пока не реализуются}

\subsection{Использование знака <<->>}

Знак <<\verb|-|>> используется для задания бинарных и унарных операций. В случае бинарных операций его аргументом является двухэлементный кортеж. Интерпретация знака <<\verb|-|>> в качестве функции зависит от типов элементов кортежа. Если аргумент является двухэлементным кортежем числовых атомов (целых или действительных), то выполняется арифметическое вычитание. При вычитании двух целых чисел результат всегда является целым числом. В этом случае автоматическое преобразование к действительному числу не происходит даже при переполнении. Вычитание действительных чисел осуществляется, когда аргументом функции является двухэлементный кортеж действительных чисел. Результатом в этом случае является действительное число. Следует отметить отсутствие неявных преобразований при выполнении вычитания. Исключением является преобразование целочисленной константы "0" в действительный нуль "0.0".

Если аргумент является числовым атомом, то в качестве результата осуществляется смена знака числа (унарный минус).

Аргумент функции <<\verb|-|>> может также быть булевым кортежем длиной, равной 2. Результатом в этом случае является исключающее или. Допускается одноэлементный булев атом, порождающий в качестве результата инверсное значение.

Во всех остальных случаях результатом является ошибка операции интерпретации \verb|BASEFUNCERROR|.

Функция, имеет следующие сигнатуры:
\begin{verbatim}
    :- << func (int, int) -> int
    :- << func (double, double) -> double
    :- << func (bool, bool) -> bool
    :- << func int -> int
    :- << func double -> double
    :- << func bool -> bool
\end{verbatim}

Примеры.
\begin{verbatim}
(3,5):- => -2
(3.0, -5.0):- => -2.0
(3, 5.0):- => BASEFUNCERROR
(5):- => BASEFUNCERROR
5:- => -5
(min_integer,-1):- => 0
(true, true):- => false
true:- => false
false:- => true
(true):- => BASEFUNCERROR
\end{verbatim}

\subsection{Использование знака <<*>>}

Знак <<\verb|*|>> используется для задания бинарной функции умножения. Его аргументом является двухэлементный кортеж. Интерпретация знака <<\verb|*|>> в качестве функции зависит от типов элементов кортежа. Если аргумент является двухэлементным кортежем числовых атомов (целых или действительных), то выполняется арифметическое умножение. При умножении двух целых чисел результат всегда является целым числом. В этом случае автоматическое преобразование к действительному числу не происходит даже при переполнении. Умножение действительных чисел осуществляется, когда аргументом функции является двухэлементный кортеж действительных чисел. Результатом в этом случае является действительное число. Следует отметить отсутствие неявных преобразований при выполнении умножения. Исключением является преобразование целочисленной константы "0" в действительный нуль "0.0".

Аргумент функции <<\verb|*|>> может также быть булевым кортежем длиной, равной 2. Результатом в этом случае является конъюнкция (логическое произведение).

Во всех остальных случаях результатом является ошибка операции интерпретации \verb|BASEFUNCERROR|.

Функция, имеет следующие сигнатуры:
\begin{verbatim}
    :* << func (int, int) -> int
    :* << func (double, double) -> double
    :* << func (bool, bool) -> bool
\end{verbatim}

Примеры.
\begin{verbatim}
    (3, 5):* => 15
    (3, 5.0):* => BASEFUNCERROR
    (3.0, 5.0):* => 15
    (5):* => BASEFUNCERROR
    5:* => BASEFUNCERROR
    (true, false):* => false
    (true, true):* => true
    true:* => BASEFUNCERROR
    (true):* => BASEFUNCERROR
\end{verbatim}

\subsubsection{Использование знака <</>>}

Знак <<\verb|/|>> используется для задания бинарной операции деления. Его аргументом является двухэлементный кортеж. Интерпретация знака <</>> в качестве функции зависит от типов элементов кортежа. Аргумент должен является двухэлементным кортежем числовых атомов (целых или действительных), что определяет операцию деления. При делении двух целых чисел результат всегда является целым числом. В этом случае автоматическое преобразование к действительному числу не происходит даже при наличии остатка. Деление действительных чисел осуществляется, когда аргументом функции является двухэлементный кортеж действительных чисел. Результатом в этом случае является действительное число. Следует отметить отсутствие неявных преобразований при выполнении деления. Исключением является преобразование целочисленной константы "0" в действительный нуль "0.0".

При делении на ноль выдается ошибка \verb|ZERODIVERROR|. Во всех остальных случаях результатом является ошибка операции интерпретации \verb|BASEFUNCERROR|.

Функция, имеет следующие сигнатуры:
\begin{verbatim}
    :/ << func (int, int) -> int
    :/ << func (double, double) -> double
\end{verbatim}

Примеры.
\begin{verbatim}
    (3, 5):/ => 0
    (3.0, 5.0):/ => 0.66667
    (3, 5.0):/ => BASEFUNCERROR
    (3.0, 0):/ => ZERODIVERROR
    (5):/ => BASEFUNCERROR
    5:/ => BASEFUNCERROR
\end{verbatim}

\subsubsection{Использование знака <<\%>>}

Знак <<\verb|%|>> используется для задания бинарной функции вычисления остатка от целочисленного деления. Его аргументом является двухэлементный целочисленный кортеж. Результат всегда является целым числом.

При втором аргументе, равном нулю, выдается ошибка \verb|ZERODIVERROR|. Во всех остальных случаях результатом является ошибка операции интерпретации \verb|BASEFUNCERROR|.

Функция, имеет следующие сигнатуры:
\begin{verbatim}
    :% << func (int, int) -> int
\end{verbatim}

Примеры.
\begin{verbatim}
    (13,5): % => 3
    (13,-5): % => 3
    (-13,5): % => -3
    (-13,-5): % => -3
\end{verbatim}

\subsection{Использование знака <</\%>>}

Знак <<\verb|/%|>> используется для одновременного выполнения двух бинарных функций: целочисленного деления и  вычисления остатка. Его аргументом является двухэлементный целочисленный кортеж. Результат является двухэлементным целочисленным кортежем, первый элемент которого содержит частное от деления, а второй - остаток.

При втором аргументе, равном нулю, выдается ошибка \verb|ZERODIVERROR|. Во всех остальных случаях результатом является ошибка операции интерпретации \verb|BASEFUNCERROR|.

Функция, имеет следующие сигнатуры:
\begin{verbatim}
    :/% << func (int, int) -> (int, int)
\end{verbatim}

Примеры.
\begin{verbatim}
    (13,5): /% => (2, 3)
    (13,-5): /% => (-2, 3)
    (-13,5): /% => (-2, -3)
    (-13,-5): /% => (2, -3)
\end{verbatim}

Подобная трактовка отличается от математической, рассматриваемой, например в первом томе книги Дональда Кнута:
\begin{verbatim}
    x = (x DIV y) * y + (x MOD y)
    0 <= (x MOD y) < y, если x > 0 или
    0 >= (x MOD y) > y, если x < 0
\end{verbatim}
Кстати, Кнут и не называет операцию MOD остатком. Получить математические версии функции, использую предопределенную операцию можно следующим образом:
\begin{verbatim}
    // функция целочисленного деления:
    div << func x@(int, int) -> int
    {
        dm << x:%:1;
        {(dm,0):(>=,<):?}^
        (
            dm,
            {(dm,1):-}
        ):return
    }

    // функция выделения остатка от целочисленного деления:
    mod << func x@(int, int) -> int
    {
        dm<<x:%;
        {(dm:1,0):(>=,<):?}^
        (
            dm:2,
            {(dm:2,0):(>=,<):?}^
            (
                {(dm:2,1):-:-},
                {(dm:2,1):+:-}
            )
        ):return
    }
\end{verbatim}


\textbf{Примечание.}
\textit{\textbf{Примечание.}
Следует отметить, что в этих выражениях обнаружен косяк, который надо бы исправить!!!
Функции нужно внимательно проверить, так как написал их без особого контроля.}


\subsection{Использование знаков <<=>>, <<!=>>, << < >>, << <= >>, << > >>, << >= >>}

Данные знаки используется для задания бинарных операции сравнения. Аргументом любого из них является двухэлементный кортеж, данные которого имеют одинаковый тип. Неявное автоматическое преобразование данных к одному типу внутри кортежа не происходит. Допускается сравнение между собой:
\begin{itemize}
    \item целочисленных данных --- все операции;
    \item действительных данных --- все операции;
    \item символов --- все операции;
    \item булевских данных --- все операции;
    \item данных, для которых тип не определен (пустой тип) --- на равенство и неравенство;
    \item типов --- на равенство и неравенство;
    \item функций --- на равенство и неравенство.
\end{itemize}
Результатом является булевское значение.
Во всех остальных случаях результатом является ошибка операции интерпретации \verb|BASEFUNCERROR|.

\textit{\textbf{Примечание.}
Появилась дурная мысль: добавить сравнение на все операции для функций. Чтобы, например, определять отношение: обобщение --- специализация. Но пока это только мысль...}

Пусть \verb|cmp| означает один из представленных выше знаков операции сравнения. Тогда функции имеют следующие обобщенны сигнатуры:
\begin{verbatim}
    :cmp << func (int, int) -> bool
    :cmp << func (double, double) -> bool
    :cmp << func (char, char) -> bool
    :cmp << func (bool, bool) -> bool
    :cmp << func (empty, empty) -> bool
    :cmp << func (type, type) -> bool
    :cmp << func (func, func) -> bool
\end{verbatim}
При этом необходимо учесть что не для каждого из представленных типов имеются все реализации операций сравнения (описано выше).

\subsection{Использование в качестве функций предопределенных типов}

Применение статической типизации позволяет в большинстве случаев знать типы артефактов во время компиляции. Также во время компиляции многие типы можно вывести из функций и других типов. Вместе с тем, в случае обобщений (при использовании типа \verb|union|) необходимо определять тип артефакта во время выполнения. Поэтому наличие функций, обеспечивающих анализ типов в целом необходимо.

Непосредственно типы определяются как перечисления. При этом предопределенные типы заданы изначально. Пользовательские типы добавляются в процессе написания программы. Все они относятся к типу типов \verb|type|. Непосредственно в функциях типы являются частью обрабатываемых значений, каждое из которых определяется в виде двойки:
\begin{verbatim}
    Значение = (тип, величина).
\end{verbatim}

Наряду с обработкой данных, осуществляемой неявно, допускается выделять тип любого элемента данных. Для этого используется предопределенная операция \verb|type|. Формируемая при этом величина принадлежит к <<типовым>> и имеет точно такую же организацию, как и любой другой аргумент. Ее специфика проявляется лишь в том, что типом аргумента является \verb|type|.
\begin{verbatim}
    Структура типового элемента = (type, значение типа).
\end{verbatim}

Поэтому для выделения типа ключевое слов \verb|type| используется в качестве функции в операторе интерпретации. В результате возвращается значение типа, принадлежащее множеству допустимых типов артефактов. Функция получение типа значения имеет следующую сигнатуру:
\begin{verbatim}
    type << func any -> type
\end{verbatim}

Примеры.
\begin{verbatim}
    10:type => int
    3.14:type => double
    (1, 2, 4, 7): type => tuple
    [1, 2, 3, 3, 4]: type => swarm
\end{verbatim}

Применение функции \verb|type| к <<типовому>> элементу невозможно и ведет к ошибке интерпретации \verb|TYPEERROR|, например:
\begin{verbatim}
    int : type => TYPEERROR
\end{verbatim}

Предопределенные типы в качестве функций могут использоваться для различных операций преобразования типов, определяемых в языке. При этом можно считать, что выступая в качестве функций они имеют определенные сигнатуры, в которых аргументы--данные, принадлежащие входному типу, преобразуются в значения, тип которого определяется результатом.

Существует также возможность, за счет использования пользовательской перегрузки, формировать собственные преобразователи как для предопределенных, так и порожденных именованных типов. В роли такого преобразователя выступает специальный оператор \verb|transform|. По сути это некоторая функция--конструктор, синтаксис которой отличается от обычной функции в связи со спецификой использования. Синтаксис преобразователь задается следующим образом:
\begin{verbatim}
    преобразователь = "transform" имя_аргумента "@" имя_типа "->" имя_типа
        [тело_преобразованиеля ].
    тело_преобразователя = "{" [ Элемент {";" Элемент }] "}".
\end{verbatim}
Практически тело преобразователя совпадает с телом функции. В результате при компиляции формируется описание сигнатуры (прототипа), которое полностью совпадает с описанием функции, имя которой совпадает с именем типа, что позволяет компилятору работать с преобразователями также как и с функциями.

\textit{\textbf{Примечание.}}
\textit{В данной ситуации при в общем-то понятной и достаточно прозрачной семантике встал вопрос, каким образом должен выглядеть синтаксис преобразователя типа. Рассматривались следующие варианты.}

\textit{1) Описание всех преобразователей в виде специальной функции трансформации, перегружаемой на все возможные ситуации. Ее схематичный вид должен был выглядеть следующим образом:}
\begin{verbatim}
    "transform" << "func" имя_аргумента "@" имя_входного_типа
            "->" имя_выходного_типа тело
\end{verbatim}
\textit{Например:}
\begin{verbatim}
    transform << func x@double -> int {...}
\end{verbatim}
Не понравилось. Появляются одинаковы сигнатуры аргументов для разных выходных типов, что не соответствует данному ранее условию использования перегружаемых функций. Смотрится не очень логично и красиво. При выводе типов нужно учитывать имя функции в коде, а затем проверять, что на выходе.

2) Другой вариант связан с непосредственной перегрузкой имени типа в качестве имени функции. Например:
\begin{verbatim}
    int << func x@double -> int {...}
\end{verbatim}
В этой ситуации логика использования сохраняется. Однако подобная реализация допускает при написании использовать разные имена функции и типа возвращаемого значения. Теоретически это возможно, но концептуально выглядит не очень хорошо. Можно для типов ввести дополнительных контроль на этапе компиляции, но в целом все это выглядит не очень красиво. Дважды повторяется имя типа.

3) Поэтому выбран вариант с вводом специального оператора \verb|transform|, что позволяет избавиться от избыточности и акцентировать внимание не на сигнатуре, а на выходном параметре. Также видна логика преобразователя, как конструктора с аргументом на входе. Для внутреннего представления и формируемого экспорта сохранено описание преобразователя как прототипа функции, определяющего имя, аргумент и результат.

\subsubsection{Использование типа <<int>>}

Функция осуществляет преобразование аргумента к целочисленному типу. Эти преобразования возможны для действительных чисел, символов, булевских величин. Действительные числа преобразуются к целым путем отсечения дробной части независимо от знака числа.

Функция для экспорта имеет следующие сигнатуры:
\begin{verbatim}
    int << func bool -> int
    int << func double -> int
    int << func char -> int
\end{verbatim}

Примеры.
\begin{verbatim}
    13.8:int  => 13
    -13.2:int  => -13
    -13.8:int  => -13
    '0':int => 48
    true:int => 1
    false:int => 0
\end{verbatim}

\subsubsection{Использование типа <<bool>>}

Функция осуществляет преобразование аргумента к булевскому типу. Эти преобразования возможны для целых чисел. Если число равно нулю, то результатом является \verb|false|. В противном случае результат равен \verb|true|.

Функция имеет следующую сигнатуру:
\begin{verbatim}
    bool << func int -> bool
\end{verbatim}

Примеры.
\begin{verbatim}
    1:bool  => true
    -10:bool  => true
    100:bool => true
    0:bool => false
\end{verbatim}

\subsubsection{Использование типа <<signal>>}

Функция осуществляет преобразование аргумента любого типа к сигнальному значению. То есть идет потеря величины аргумента.
Функция имеет следующую сигнатуру:
\begin{verbatim}
    signal << func any -> signal
\end{verbatim}

Примеры.
\begin{verbatim}
    1:signal  => !
    !:signal  => !
\end{verbatim}

\subsubsection{Использование типа <<error>>}

Функция осуществляет формирование общей ошибки и прерывание программы. Может применяться к любому аргументу.

Функция имеет следующую сигнатуру::
\begin{verbatim}
    error << func any -> ERROR
\end{verbatim}

Примеры::
\begin{verbatim}
    1:error  => ERROR
    !:error  => ERROR
\end{verbatim}

\subsection{Использование в качестве функций имен составных типов}

Формируемые пользователем типы в основном базируются на составных типах, позволяющих конструировать разнообразные конфигурации данных, которым присваиваются соответствующие имена. Эти имена пользовательских типов также могут применяться для задания пользовательских функций. Однако в ряде преобразований как уже сформированных пользовательских типов, так и программоформирующих операторов может осуществляться путем использования имен, определяющих составные типы. К этим именам относятся: \verb|vector|, \verb|tuple|, \verb|struct|, \verb|swarm|, \verb|stream|, \verb|union|.

Для поддержки строгой типизации в языке запрещены неявные преобразования. С одной стороны это ведет к дополнительным проблемам для программиста. Но с другой стороны повышается надежность кода за счет усиления контроля над типами.

\subsubsection{Использование типа <<vector>>}

Задает преобразование аргумента в вектор. Аргумент может быть базовым элементарным типом. В этом случае он становится одноэлементным вектором. Также эта функция может использоваться с роями, преобразуя их в вектор. Это позволяет синхронизировать поступающие данные перед последующей обработкой.
\begin{verbatim}
    атом:vector => (атом)
    [элемент, ... элемент]:vector => vector(элемент, ... элемент)
\end{verbatim}

Функция также может применяться к потоку, который в этом случае преобразуется к вектору. Тем самым поэлементный вывод из потока трансформируется в формирование вектора, число элементов в котором определяется количеством элементов, порождаемых потоком (что не всегда хорошо).
\begin{verbatim}
    stream{}:vector => vector(элемент, ... элемент)
\end{verbatim}

Данная функция имеет следующие сигнатуры
\begin{verbatim}
    vector << func any -> any(*)
    vector << func any[*] -> any(*)
    vector << func any{*} -> any(*)
\end{verbatim}

Помимо этого можно задавать ряд преобразований, если они возможны во время компиляции. В частности допускается преобразовывать структуру, состоящую из элементов одного типа в вектор. Аналогичное преобразование в вектор возможно и для кортежа, все элементы которого имеют один тип. Для этих преобразований имеются соответствующие функции::
\begin{verbatim}
vector << func tuple([n]@any) -> any(n)
vector << func struct([n]@any) -> any(n)
\end{verbatim}

Через \verb|[n]@any| обозначено наличие в кортеже или структуре \verb|n| элементов одного любого типа. Если аргумент не соответствует сигнатуре, то результатом является ошибка операции интерпретации \verb|BASEFUNCERROR|.

\textit{\textbf{Примечание.}
Здесь и ниже нужно сделать коррекцию в описаниях, связанную с тем, что произошло определенное разделение между типом контейнера и типом активного объекта. Нужно показать, что в данном случае идет преобразование типа контейнера при сохранении типа или вида объекта. Пока же наблюдается помесь, обусловленная тем, что ранее это разделение не просматривалось. Коррекцию сделать при следующем прочтении, так как на текущем этапе пока идет перенос текста в тех.}

\subsubsection{Использование типа <<tuple>>}

Задает преобразование аргумента в кортеж. Аргумент может быть базовым элементарным типом. В этом случае он становится одноэлементным кортежем. Также эта функция может использоваться с соединителям и роями, сохраняя свойственные для этих активных объектов механизм обработки данных.
\begin{verbatim}
    атом:tuple => tuple(атом)
    [элемент, ... элемент]:tuple => tuple(элемент, ... элемент)
\end{verbatim}

Функция также может применяться к потоку (???), который в этом случае преобразуется к кортежу. Тем самым поэлементный вывод из потока трансформируется в формирование кортежа, число элементов в котором определяется количеством элементов, порождаемых потоком (что не всегда хорошо).
\begin{verbatim}
    stream{}:tuple => tuple(элемент, ... элемент)
\end{verbatim}

Данная функция имеет следующие сигнатуры:
\begin{verbatim}
    tuple << func any -> tuple(any)
    tuple << func any(*) -> tuple([n]@any)
    tuple << func any[*] -> tuple([n]@any)
\end{verbatim}

Другим аналогичным преобразованием является преобразование в кортеж структуры. В этом случае именованный доступ к полученному результату заменяется на позиционный. Порядок следования типов данных в кортеже совпадает с порядком их следования в исходной структуре. Преобразование описывается следующей прототипом:
\begin{verbatim}
    tuple << func struct(a1:t1, a2:t2, ...) -> tuple(t1, t2, ...)
\end{verbatim}

Если аргумент не соответствует сигнатуре, то результатом является ошибка операции интерпретации \verb|BASEFUNCERROR|.

\subsubsection{Использование типа <<swarm>> и спецсимвола <<[]>>}

Задает преобразование поступающего активного объекта к рою с сохранением его контейнерного типа. Аргумент может быть базовым элементарным типом. В этом случае он становится одноэлементным вектором. Функция может использоваться с различными контейнерными типами. Ее применении к соединителю позволяет сделать поступающие данные асинхронными, что проявляется в ходе последующей их обработки. Компилятор подобные преобразования может учесть в момент анализа программы.
\begin{verbatim}
    атом:swarm => [атом]
    (элемент, ... элемент):swarm => [элемент, ... элемент]
\end{verbatim}

Функция также может применяться к потоку, который в этом случае преобразуется к рою. Тем самым поэлементный вывод из потока трансформируется в формирование роя, число элементов в котором определяется количеством элементов, порождаемых потоком. Порядок элементов в рое совпадает с последовательностью их порождения потоком. Информация о каждом элементе поступает в оператор интерпретации асинхронно.
\begin{verbatim}
    stream{}:swarm => [элемент, ... элемент]
\end{verbatim}

Данная функция имеет следующие сигнатуры::
\begin{verbatim}
    swarm << func any -> [any]
    swarm << func any(*) -> any[*]
    swarm << func any{*} -> any[*]
\end{verbatim}
Если аргумент не соответствует сигнатуре, то результатом является ошибка операции интерпретации \verb|BASEFUNCERROR|.

\subsubsection{Использование типа <<stream>>}

Задает преобразование к потоку. Аргумент может быть базовым элементарным типом. В этом случае он становится потоком, порождающим один элемент. Также эта функция может использоваться с векторами соединителей и роев, трансформируя их в поток. Это позволяет сделать данные последовательно асинхронными, поступающими в произвольном порядке, так как компилятор подобные преобразования может учесть в момент анализа программы. То есть преобразование может осуществиться до окончательного вычисления вектора или вообще до начала его вычисления.
\begin{verbatim}
    атом:stream => stream{атом}
    vector(элемент, ... элемент):stream => stream{элемент, ... элемент}
\end{verbatim}

Данная функция имеет следующие сигнатуры:
\begin{verbatim}
    stream << func any -> any{*}
    stream << func any(*) -> any{*}
    stream << func any{*} -> any{*}
\end{verbatim}
Если аргумент не соответствует сигнатуре, то результатом является ошибка операции интерпретации \verb|BASEFUNCERROR|.

\subsection{Использование в качестве функций пользовательских типов}

Разработчик программного обеспечения может создавать свои описания типов на основе уже существующих предопределенные типов и ранее разработанных именованных пользовательских типов. Как и предопределенные типы, именованные пользовательские типы могут использоваться в качестве функций. Однако алгоритмы, выполняемые этими функциями, а также их сигнатуры определяются явно разработчиками ПО с использованием преобразователя \verb|transform|.

Такой подход говорит о том, что в принципе пользователь может связать с такой функцией какие угодно вычисления, тем самым обеспечивая гибкость в использовании имен типов в качестве функций. Поэтому невозможно заранее предугадать их поведение.

\subsection{Использование в качестве функции обобщенного типа <<union>>}

Обобщенный глобальный тип \verb|union| в качестве признака фиксирует имя включаемого в него типа. Функция с именем \verb|union|, формирует вокруг конкретного типа оболочку, включающую признак (тэг) и ссылку на охватываемый артефакт. В качестве охватываемого типа может использоваться любой именованный тип за исключением самого \verb|union|. Таким образом, преобразование типа заключается в его охвате, как и в случае с пользовательским обобщением.

Использование в качестве функции описывается следующим прототипом:
\begin{verbatim}
    union << func имя_типа -> union
\end{verbatim}

Для выделения конкретного артефакта из глобального обобщения \verb|any| необходимо проверить признак, после чего выделить значения используя имя полученного типа в качестве функции, осуществляя тем самым явное приведение вложенного артефакта к его типу. Проверка признака осуществляется функцией `\verb|tag|, которая применима ко всем обобщениям. Если тип вложенного артефакта совпадает с типом, используемым в качестве функции, то результатом является артефакт указанного типа. В противном случае генерируется ошибка \verb|TYPEERROR| и происходит прерывание программы.

\textit{\textbf{Примечание.}
Нужен пример.}

\section{Семантика хранилищ}

Наряду с тем, что различные типы данных используются в качестве аргументов в операторе интерпретации, они также имеют конкретную структуру, которая определяет специфику представления этих данных в хранилищах. Особенностью хранилищ является то, что к ним применяется ограниченный набор функций. Этот набор ограничен чтением из хранилищ и записью в них (с соблюдением принципа единственного присваивания). Но даже описание операций чтения--записи требует определенного синтаксического сахара для их более удобного и понятного представления в языке.

\subsection{Хранилища атомарных данных}

Чтение атомарных данных задается обращением к хранилищам как к аргументам по имени хранилища. То есть, здесь никаких сложностей не возникает. Обращение обычно эквивалентно использованию переменных в традиционных языках программирования.

Пример::
\begin{verbatim}
    // чтение из хранилищ a, b и сложение
    (a, b):+
\end{verbatim}

Запись в атомарное хранилище также осуществляется по аналогии с записью в обычные переменные. Однако особенностью в данном случае является то, что хранилище используется в качестве функции в операторе интерпретации, который принимает поступающие данные и записывает их в хранилище::
\begin{verbatim}
// запись данных в целочисленное хранилище a
a@int
10:a
\end{verbatim}
Следует отметить, что тип записываемых данных должен совпадать с типом хранилища.

\subsection{Общие особенности хранилищ составного типа}

Независимо от типа составного хранилища, чтение и запись в них осуществляется одинаково в случае, когда каждое из таких хранилищ построено на основе именованного типа. В этом случае описание хранилища практически не отличается от описания хранилищ атомарных типов. Чтение осуществляется обращением целиком ко всем данным, а селективная выборка при этом может определяться соответствием между типом хранилища и функцией, задаваемой в операторе интерпретации. При этом допускается неявное приведение типов данных в соответствии со сформированной сигнатурой функции.

Простейший вариант чтения заключается во взаимодействии одноименных по типу хранилищ, когда данные из одного копируются в другое::
\begin{verbatim}
    intPair << @(int, int)
    a:intPair
    b:intPair
    (3,4):intPair:a     // явное преобразование типа tuple (int, int)
    a:b                 // запись из a в b где b рассматривается как функция
    // с сигнатурой: b << func x@Pair -> signal
\end{verbatim}
Вместе с тем следует отметить, что специфические особенности организации каждого из хранилищ ведут к ряду особенностей, связанных с обращением к ним.

\subsection{Особенности хранилищ векторного типа}

Наряду с непосредственным копированием данных между векторами одинакового типа и размера возможны различные варианты организации чтения из вектора элемента по индексу. Первый вариант связан с использованием целочисленных величин в качестве индексов. Он позволяет прочитать элемент в качестве результата::
\begin{verbatim}
    (1, 2, 3, 4, 5):vector: 3 => 3
    (1, 2, 3, 4, 5):vector: i => результат - i-й элемент вектора
\end{verbatim}


Другой вариант связан с представлением операции обращения к вектору с использованием специального его обозначения:
\begin{verbatim}
    ЭлементВектора = ИмяВектора "(" Индекс ")".
\end{verbatim}


Данный вариант полностью эквивалентен предыдущему, но может использоваться только в случае, когда вектор имеет имя:
\begin{verbatim}
    A << (1, 2, 3, 4, 5):vector
    A(3) => 3
    B << (1, 2, 3, 4, 5):vector
    B(i) => результат - i-й элемент вектора
\end{verbatim}
В качестве индекса допускается как целочисленная константа, так и целочисленное значение, полученное в ходе вычислений. Следует отметить, что в первом случае выход за границы вектора определяется компилятором.

Запись в вектор возможна только в том случае, если он предварительно объявлен как незаполненное хранилище. При этом конструкция, определяющая элемент вектора используется как функция в операторе интерпретации:
\begin{verbatim}
    A@int(5)
    7:A(i)
    10:A(1)
\end{verbatim}
При попытке повторно записать в уже заполненный элемент, запрещенное принципом единственного присваивания, происходит прерывание программы

Особенности хранилищ имеющих тип кортеж
~~~~~~~~~~~~~~~~~~~~~~~~~~~~~~~~~~~~~~~~~~~~~~~~~~~~~~~~~~~~~~~~~~~~~~~~~~~~~~

Чтение-запись с использованием кортежей во многом напоминает аналогичные операции с векторами. Однако имеются и свои особенности. Спецификой кортежей является наличие в них элементов разных типов. Поэтому основной вид доступа связан с использованием в качестве индекса целочисленной константы. В этом случае компилятор может точно определить тип результата::
\begin{verbatim}
    (1, 2, 3, 4, 5): 3 => 3@int
    (3, true, !): 2 => true@bool
\end{verbatim}

Ошибка компиляции происходит в том случае, если в качестве индекса выступает величина вычисляемая во время выполнения программы. Однако имеется исключение. Если все элементы кортежа имеют один и тот же тип, компилятор может допустить использование вычисляемой величины, так как тип результата становится известным::
\begin{verbatim}
    (1, 2, 3, 4, 5): i =>
            результат - i-й элемент кортежа однотипных элементов типа int
    (3, true, !): j => ошибка компиляции
\end{verbatim}

Для организации чтения может использоваться целое число как функция оператора интерпретации или обозначение элемента кортежа, аналогичное по синтаксису элементу вектора:
\begin{verbatim}
    ЭлементКортежа = ИмяКортежа "(" Индекс ")".
\end{verbatim}

Запись в кортеж-хранилище осуществляет по тем же правилам, что и чтение. В общем случае она происходит, если индекс задается константой, позволяющий сопоставить типы записываемых данных и поля кортежа. Вычисляемая величина допускается в качестве индекса, если все поля кортежа однотипны.
\begin{verbatim}
    X@(int, bool, signal)
    false:X(2)
    10:A(1)
\end{verbatim}

\subsection{Особенности структурных хранилищ}

По разнотипности своих элементов структуры аналогичны кортежам. Однако для доступа к полям структуры используются внутренние имена этих полей, с каждым из которых сопоставлен тип. Поэтому чтение-запись реализованы только с использованием этих полей без каких либо вычислений. Обычный доступ осуществляется через имя хранилища, после которого через точку следует имя поля:
\begin{verbatim}
    ЭлементСтруктуры = ИмяХранилища "." ИмяПоля.
\end{verbatim}

Однако возможны ситуации, когда результат структурного типа формируется в ходе вычислений. В этом случае имя поля, предваряемого точкой, может быть указано в качестве функции в операторе интерпретации:
\begin{verbatim}
    TS << @(x@int, b@bool)
    S@TS
    // Здесь какие-то вычисления порождающие структуру в качестве результата
    ...:TS:.b  // После приведения к типу TS на выходе будет значение поля b
    S.a => возвращает значение поля a структуры S
    S:.a => аналогичный результат
\end{verbatim}

Запись в структуру-хранилище осуществляется с явным использованием элемента структуры в качестве аргумента-функции оператора интерпретации:
\begin{verbatim}
    S@(i@int, b@bool, s@signal)
    false:S.b
    10:S.i
\end{verbatim}

\subsection{Особенности хранилищ роевого типа}

Спецификой роя является асинхронное поступление данных и, как следствие отдельная сигнализация о поступлении каждого элемента в связанный с ним оператор интерпретации. Вместе с тем обращение к рою по индексу практически аналогично обращению к вектору. При отсутствии элемента происходит ожидание его появления. Поэтому возможны те же варианты:
\begin{verbatim}
    [1, 2, 3, 4, 5]: 3 => 3
    [1, 2, 3, 4, 5]: i => результат - i-й элемент роя
\end{verbatim}

Другой вариант связан с представлением операции обращения к рою с использованием специального его обозначения:
\begin{verbatim}
    ЭлементРоя = ИмяРоя "[" Индекс "]".
\end{verbatim}

Данный вариант полностью эквивалентен предыдущему, но может использоваться только в случае, когда Рой имеет имя:
\begin{verbatim}
    R << [1, 2, 3, 4, 5]
    A[3] => 3
    Q << [1, 2, 3, 4, 5]
    Q[i] => результат - i-й элемент роя
\end{verbatim}
В качестве индекса допускается как целочисленная константа, так и целочисленное значение, полученное в ходе вычислений. Следует отметить, что в первом случае выход за границы роя определяется компилятором.

Запись в рой возможна только в том случае, если он предварительно объявлен как незаполненное хранилище. При этом конструкция, определяющая элемент роя используется как функция в операторе интерпретации:
\begin{verbatim}
    R@int[5]
    7:R[i]
    10:R[1]
\end{verbatim}
При попытке повторно записать в уже заполненный элемент, запрещенное принципом единственного присваивания, происходит прерывание программы

\subsubsection{Реагирование и доступ по мере поступления элементов в рой}

Если рой является аргументом, то оператор интерпретации реагирует на появление каждого элемента. В этом случае можно использовать функции, которые также срабатывают на каждый элемент и обрабатывают рой по мере поступления отдельных данных, не дожидаясь полной синхронизации, как это происходит в случае вектора, кортежа или структуры.

Для реализации дополнительных возможностей предполагается использовать ссылки, каждая из которых имеет дополнительные независимые механизмы по контролю за данными поступающими в рой.

\textit{\textbf{Примечание.}
Данный фрагмент нужно еще обдумать и дописать}

\subsection{Особенности хранилищ - потоков}

Основной спецификой потоков является отсутствие информации о количестве поступающих данных. В результате завершение очереди определяется по достижению "маркера дна". То есть, признака, указывающего на то, что данные в потоке отсутствуют и больше не появятся. В связи с этим возникает определенная специфика в формировании семантики этого артефакта.

В целом поведение очереди во многом аналогично поведению асинхронного списка, реализованного в языке программирования Пифагор. Однако ряд специфических моментов требуют других решений при манипуляции с ним.

Первая особенность заключается в том, что при чтении данных может возвратиться или значение, или признак завершения данных. При использовании статической типизации это два разных вида информации, поэтому необходимо разделение возвращаемых данных, которое можно сделать реализовав две функции:
\begin{itemize}
    \item функцию проверки на наличие в потоке данных (\verb|is|), которая возвращает \verb|true|, если данные еще поступают, и \verb|false|, если данные в поток поступать перестали;
    \item функцию получения элемента из потока (\verb|get|), которая при наличии в нем данных выбирает первый элемент, тип которого определяется из описания потока, а отсутствие данных ведет к прерыванию программы.
\end{itemize}

Помимо этого необходимо вытолкнуть из потока прочитанный элемент, чтобы иметь доступ к следующему. Для этого можно использовать функцию \verb|pop|.

Вместе с тем, использование нескольких функций может привести к коллизиям, если к потоку идет одновременное обращение из нескольких точек программы. Для их предотвращения необходимо использовать атомарные функции, выполняющие все необходимые действия за один раз таким образом, чтобы при каждом обращении происходила корректная обработка данных. Можно например, совместить проверку с чтением элемента, возвращая пару значений: (данные, признак наличия данных). Если признак наличия данных будет равен \verb|true|, то прочитаны корректные данные. Их можно обрабатывать. В противном случае данные неактуальны, так как поток завершил их порождение. Однако в этом случае коллизии могут все равно возникать и быть связаны с выталкиванием данных.

Более надежным смотрится вариант, когда в одной атомарной функции совмещаются все три, описанные выше. То есть, когда при наличии данных происходит их выдача в виде двойки и выталкивание элемента из потока. Тогда каждая из операций чтения будет брать только свои уникальные данные и убирать их из потока. При обращении множества операций к пустому потоку все они получат значение признака наличия данных, равное \verb|false|. При этом прерывания программы не происходит. Предполагается, что в качестве такой неделимой операции будет использоваться применением очереди в качестве данных оператора интерпретации. Например:
\begin{verbatim}
    S@int{}
    (S:1, S:1):+
\end{verbatim}

В примере демонстрируется двойное обращение к потоку, после чего полученные числа складываются. В данной ситуации порядок выдачи данных из очереди недетерминирован. То есть неясно, в какой аргумент функции сложения попадет из потока первый элемент данных. Для реализации детерминированности необходимо формировать дополнительные сигналы обеспечивающие требуемую последовательность:
\begin{verbatim}
    S@int{}
    x1 << S:! (??? может быть обращение задавать сигналом?)
    (x1, S:{x1:signal}):+
\end{verbatim}


\textit{\textbf{Примечание.}
Расширить возможности манипуляций с потоками, поддерживающими многократное дублирование доступа, можно через ссылки. Но это отдельная тема для анализа и реализации. Пока нет.}

Запись в поток осуществляется за счет использования его имени в качестве функции:
\begin{verbatim}
    S@int{}
    13:S
\end{verbatim}

Можно также осуществлять непосредственную передачу данных из одного потока в другой, которая может сочетаться с передачей независимых данных. Порядок передачи при этом также может быть недетерминированным:
\begin{verbatim}
    S1@int{}
    S2@int{}
    S1:S2
    5:S2
    10:S1
\end{verbatim}


\textit{\textbf{Примечание.}
Более детальную проработку описания очередей предполагается сделать, когда начнется разработка и анализ реализация функций - примеров.}

\section{Семантика объектов данных}

Объекты данных обеспечивают группировку поступающих величин, используя предопределенные методы. Они также могут использоваться в качестве данных и функций в операторе интерпретации. Они имеют те же типы, что и хранилища. Основное отличие от хранилищ заключается в отсутствии предварительного описания. Поэтому их атрибуты формируется во время их описания внутри функций. При этом использование статической типизации позволяет вывести тип во время компиляции и использовать полученный результат для контроля взаимодействия различных артефактов.

\textit{\textbf{Примечание.}
??? Следует отметить, что представленные ниже функции практически дублируют их использование в операторе групповой интерпретации. Это ведет к определенной избыточности языка, что не является хорошей манерой. Поэтому, скорее всего, использование программоформирующих операторов в качестве функций пока будет запрещено. Возможно в перспективе, если появятся отличия в реализации, этот вопрос будет снова поднят. Ниже лежащий текст подраздела пока оставлен. Может быть он будет перенесен в раздел описывающий семантику оператора групповой интерпретации. ???}

\textit{Возможно, что изменения могут затронуть именно групповую интерпретацию...}

\textit{Последнее ощущение, что в данном операторе оставлю как есть...}

\subsection{Оператор группировки в кортеж}

Используется для объединения как однотипных, так и неоднотипных данных. Если группируются данные одного типа, то результат группировки может обрабатываться как аргумент функции, манипулирующей векторами. Основная задача данного оператора заключается в синхронизации поступающих данных перед последующим использованием. То есть, готовность к выполнению оператора интерпретации, в котором кортеж используется в качестве одного из аргументов сформируется тогда, когда будут вычислены все элементы кортежа. Кортеж может являться аргументом--данными как в предопределенных функциях, так и в функциях разрабатываемых пользователем. В этом случае особенности его выполнения связаны с особенностями  семантики аргумента--функции оператора интерпретации.

С другой стороны кортеж может использоваться в качестве аргумента функции одиночного оператора интерпретации. В этой ситуации каждый элемент внутри кортежа определяется как отдельная функция, которая осуществляет обработку одного и того же аргумента-данных. Результатом выполнения является кортеж, в котором собраны все вычисленные величины. Можно привести эквивалентное преобразование, определяющее одну из формул алгебры преобразования языка программирования::

\verb|  data:(f1, f2, f3, ... fn)| $\equiv$ \verb|(data:f1, data:f2, data:f3, ... data:fn)|

Пример::
\begin{verbatim}
    (10, 3):(+, -, *, /, %) => (13, 7, 30, 3, 1)
\end{verbatim}


Оператор группировки в рой
~~~~~~~~~~~~~~~~~~~~~~~~~~~~~~~~~~~~~~~

Используется для асинхронной группировки, когда появление каждого элемента сопровождается передачей информации оператору интерпретации. Использование роя в качестве аргумента-данных зависит от того, какая функция над ним выполняется. Когда рой является аргументом-функцией оператора интерпретации, то его поведение во многом аналогично использованию в качестве функции кортежа. Вместе с тем, если функции вычисляются, то есть, формируются в разные моменты времени, то обработка данных осуществляется асинхронно, также независимо на выходе появляются результаты вычислений, которые группируются в соответствующем рое. Следующая формула описывает эквивалентное преобразование роя::

\verb|  data:[f1, f2, f3, ... fn]| $\equiv$ \verb|[data:f1, data:f2, data:f3, ... data:fn]|

Пример::
\begin{verbatim}
    (10, 3):[+, -, *, /, %] => [13, 7, 30, 3, 1]
\end{verbatim}


Следует отметит, что все элементы роя, выступающие в роли функций, по определению должны иметь одинаковую сигнатуру и порождать на выходе однотипные результаты в отличие от результатов, которые могут порождать функции, сгруппированные в кортеж.

\subsection{Оператор группировки в поток}

Данный оператор используется для формирования асинхронно порождаемых последовательностей. Как и другие операторы группировки он позволяет объединить несколько элементов окончательно формируемые значения которых должны быть одного типа. Спецификой данного оператора является отсутствие фиксированной длины, что позволяет использовать в качестве элементов другие операторы группировки в очередь, порождающие данные того же типа, что и тип оператора группировки. Использование в качестве аргумента--данных, как и в предыдущих операторах, предполагает, что семантика определяется аргументом--функцией, обрабатывающей данную поток. Семантика при использовании оператора в качестве аргумента-функции определяется теми элементами, которые поступают в поток и применяются в порядке поступления к каждому аргументу. Следует отметить, что в связи с асинхронностью и разным временем обработки аргумента-данных расположение результатов в выходном потоке может отличаться от порядка функций в исходном потоке. Обработка данных множеством функций внутри потока может быть представлена следующим образом::

\verb|  data:stream{f1, f2, f3, ... fn}| $\equiv$ \verb|stream{data:f1, data:f2, data:f3, ... data:fn}|

Пример выполнения, демонстрирующий один из вариантов возможной недетерминированной перестановки результатов на выходе, что определяется спецификой обработки::
\begin{verbatim}
    (10, 3):stream{+, -, *, /, %} => stream{7, 13, 1, 3, 30}
\end{verbatim}

\subsection{Оператор задержки}

Основное назначение данного оператора заключается в задержке вычислений описанного в нем выражения независимо от того будет результат использоваться в качестве данных или функции. Оператор интерпретации раскрывает задержку, что ведет к ранее отложенному вычислению результата, тип которого определяется в зависимости от написанного выражения. Дальнейшая обработка полученной величины оператором интерпретации зависит от этого типа. Каких-либо дополнительных особенностей семантики оператора задержки нет.

\textit{\textbf{Примечание.}
Нужен пример}

Следует отметить, что в принципе возможны и другие эквивалентные способы группировки и выполнения вычислений, аналогичные описанным. Они возникают из-за возможности предварительного описания величин с последующей записью в них данных с применением принципа единственного присваивания. Также свой вклад может внести использование массового оператора интерпретации. Все эти варианты предполагается рассмотреть и сравнить в разделе примеров, описывающих технику программирования. Также обобщение вариантов предполагается сделать в разделе по алгебре эквивалентных преобразований.

\section{Использование функций с массовым оператором интерпретации}
---------------------------------------------------------------------

Применение массового оператора интерпретации в первую очередь ориентировано на использование одной функции для обработки множества данных (эквивалентно стратегии ОКМД в классификации Флинна). Предполагается что данные одного типа размещаются в контейнере (вектор, рой, очередь) и к ним применяется одна и та же функция, формируя на выходе контейнер результатов такого же типа. Подобные преобразования можно описать следующими выражениями:
\begin{verbatim}
\end{verbatim}

\verb|  (d1, d2, ... dn)::f| $\equiv$ \verb|(d1:f, d2:f, ... dn:f)|

\verb|  [d1, d2, ... dn]::f| $\equiv$ \verb|[d1:f, d2:f, ... dn:f]|

\verb|  stream{d1, d2, ... dn}::f| $\equiv$ \verb|stream{d1:f, d2:f, ... dn:f}|

Следует отметить, что в случае очереди результаты на выходе могут не соответствовать порядку, определяемому исходными данными, так как возможны недетерминированные перестановки, связанные с разным временем выполнения функций.

Помимо этого возможна ситуация, когда один элемент данных обрабатывается нескольким функциями::

\begin{center}
$\mathbf{d::(f_0, f_1, ..., f_{k-1}) \equiv (d:f_0, d:f_1, ..., d:f_{k-1})}$\\
$\mathbf{d::[f_0, f_1, ..., f_{k-1}] \equiv [d:f_0, d:f_1, ..., d:f_{k-1}]}$\\
$\mathbf{d::<[f_0, f_1, ..., f_{k-1}]> \equiv <[d:f_0, d:f_1, ..., d:f_{k-1}]>}$
\end{center}

То есть, данная ситуация эквивалентна случаю, связанному с использованием программоформирующих операторов в качестве функций.

\textit{\textbf{Примечание.}
Возможно это тот повод, который позволит отменить использование при одиночной интерпретации программоформирующие операторы в качестве функций. Зачем дублировать? Может тогда и другие проблемы и вопросы будут сняты. Надо думать...}

\textit{Но пока хочется одиночный вариант оставить...}

Другим фактором, определяющим специфику массового оператора интерпретации является то, что аргумент--функция может также задаваться с использованием программоформирующих операторов. Поэтому в ряде случаев возможно появление множества функций, обрабатывающих множество очередей. При этом контейнерные типы данных могут отличаться от контейнерных типов функций. То есть возможны следующие комбинации:

%\centering
\begin{center}
$\mathbf{(d_0, d_1, ... d_{n-1})::(f_0, f_1, ... f_{k-1}) \equiv}$ \\
$\mathbf{\equiv (d_0:(f_0, f_1, ... f_{k-1}), d_1:(f_0, f_1, ... f_{k-1}), ... d_{n-1}:(f_0, f_1, ... f_{k-1})) \equiv}$ \\
$\mathbf{\equiv ((d_0:f_0, d_0:f_1, ... d_0:f_{k-1}), (d_1:f_0, d_1:f_1, ... d_1:f_{k-1}), ...,} $\\
$\mathbf{...,  (d_{n-1}:f_0, d_{n-1}:f_1, ... d_{n-1}:f_{k-1}))}$
\newline
\newline
$\mathbf{[d_0, d_1, ... d_{n-1}]::(f_0, f_1, ... f_{k-1}) \equiv}$ \\
$\mathbf{\equiv [d_0:(f_0, f_1, ... f_{k-1}), d_1:(f_0, f_1, ... f_{k-1}), ... d_{n-1}:(f_0, f_1, ... f_{k-1})] \equiv}$ \\
$\mathbf{\equiv [(d_0:f_0, d_0:f_1, ... d_0:f_{k-1}), (d_1:f_0, d_1:f_1, ... d_1:f_{k-1}), ...,}$\\
$\mathbf{...,  (d_{n-1}:f_0, d_{n-1}:f_1, ... d_{n-1}:f_{k-1})]}$
\newline
\newline
$\mathbf{[d_0, d_1, ... d_{n-1}]::[f_0, f_1, ... f_{k-1}] \equiv}$ \\
$\mathbf{\equiv [d_0:[f_0, f_1, ... f_{k-1}], d_1:[f_0, f_1, ... f_{k-1}], ... d_{n-1}:[f_0, f_1, ... f_{k-1}]] \equiv}$ \\
$\mathbf{\equiv [[d_0:f_0, d_0:f_1, ... d_0:f_{k-1}], [d_1:f_0, d_1:f_1, ... d_1:f_{k-1}], ...,}$\\
$\mathbf{...,  [d_{n-1}:f_0, d_{n-1}:f_1, ... d_{n-1}:f_{k-1}]]}$
\end{center}
и далее все комбинации, включая потоки.

То есть, в этом случае в начале контейнерная функция применяется к каждому из элементов данных, после чего идет дальнейшее разложение. Как и в более простых ситуациях комбинации с потоками ведут к недетерминированным вычислениям и смене порядка следования результатов относительно порядка в исходных данных.

